\documentclass[a4paper,11pt]{article}
\usepackage[verbose,a4paper,tmargin=2cm,bmargin=2cm,lmargin=2.5cm,rmargin=2.5cm]{geometry}
\usepackage[utf8]{inputenc}
\usepackage{polski}
\usepackage{amsmath}
\usepackage{amsfonts}
\usepackage{amssymb}
\usepackage{lastpage}
\usepackage{indentfirst}
\usepackage{verbatim}
\usepackage{graphicx}
\usepackage{fancyhdr}
\usepackage{listings}
\usepackage{hyperref} 
\usepackage{xcolor}
\usepackage{tikz}
\usepackage{float}
\frenchspacing
\pagestyle{fancyplain}
\fancyhf{}

\usepackage{setspace}

\renewcommand{\headrulewidth}{0pt}
\renewcommand{\footrulewidth}{0.4pt}
\newcommand{\degree}{\ensuremath{^{\circ}}} 
\fancyfoot[L]{E-Biznes: Zbigniew Nowacki, Karol Podlewski, Patrycja Szczakowska}
\fancyfoot[R]{\thepage\ / \pageref{LastPage}}


\begin{document}

\begin{titlepage}
\begin{center}
\begin{tabular}{rcl}
\begin{tabular}{|r|}
\hline \\
\large{\underline{234102~~~~~~~~~~~~~~~~~~~~~~~} }\\
\small{\textit{Numer indeksu}}\\
\large{\underline{Zbigniew Nowacki~~~~~~~~} }\\
\small{\textit{Imię i nazwisko}}\\\\ \hline
\end{tabular} 
&
\begin{tabular}{|r|}
\hline \\
\large{\underline{234106~~~~~~~~~~~~~~~~~~~~~~~} }\\
\small{\textit{Numer indeksu}}\\
\large{\underline{Karol Podlewski~~~~~~~~~~~} }\\
\small{\textit{Imię i nazwisko}}\\\\ \hline
\end{tabular} 
&
\begin{tabular}{|r|}
\hline \\
\large{\underline{234121~~~~~~~~~~~~~~~~~~~~~~~} }\\
\small{\textit{Numer indeksu}}\\
\large{\underline{Patrycja Szczakowska~~~~} }\\
\small{\textit{Imię i nazwisko}}\\\\ \hline
\end{tabular} 

\end{tabular}
~\\~\\~\\ 
\end{center}
\begin{tabular}{ll}
\LARGE{\textbf{Kierunek}}& \LARGE{Informatyka Stosowana} \\
\LARGE{\textbf{Stopień}}& \LARGE{II} \\
\LARGE{\textbf{Specjalizacja}}& \LARGE{Data Science} \\
\LARGE{\textbf{Semestr}}& \LARGE{1} \\\\
\LARGE{\textbf{Data oddania}}& \LARGE{19 marca 2020} \\\\\\\\\\\\\\
\end{tabular}

\begin{center}
\textbf{\huge{E-Biznes}}
\textbf{\LARGE{\\Etap 2: Projekt}}
\textbf{\Huge{\\~\\~\\System rezerwacji\\obiektów sportowych\\~\\ JASTRZĘBIK}} 
\end{center}

\end{titlepage}

\setcounter{page}{2}
\setstretch{1.4}

\tableofcontents
\setstretch{1.2}
\newpage

\section{Charakterystyka wybranej branży}

    \subsection{Sport w Polsce}
        Zgodnie z danymi udostępnionymi przez Główny Urząd Statystyczny (GUS) w 2016 r. w Polsce działało 14,9 tys. klubów sportowych. Ich liczba w stosunku do 2014 r. wzrosła o 6,1\%. W okresie 1994-2016 liczba klubów ogółem wzrosła ponad pięciokrotnie (począwszy od 1994 r. zaczęły w Polsce powstawać uczniowskie i wyznaniowe kluby sportowe), a liczba ćwiczących wzrosła trzykrotnie. W omawianym okresie systematycznie przybywało klubów sportowych i osób ćwiczących. 
    
        Według GUS w 2016 r. kluby sportowe liczyły ponad 1 019,6 tys. członków (wzrost o 8,3 \%w stosunku do 2014 r.), a ćwiczyło w nich ponad 1 030,2 tys. osób (wzrost o 12,1\%w stosunku do 2014 r.). Podobnie jak w latach poprzednich zdecydowaną większość ćwiczących stanowili mężczyźni – 74,8\%. Także niezmiennie wśród ćwiczących dominuje młodzież do 18 lat – 71,3\%. 

        Niezmiennie w Polsce wśród ćwiczących najpopularniejsze są sporty drużynowe (58,4\% osób ćwiczących), a spośród nich piłka nożna, którą uprawia 39,5\% ćwiczących (łącznie z piłką nożną plażową i piłką nożną halową – 40,8\%). Pod względem liczby ćwiczących drugie miejsce zajmuje piłka siatkowa - 7,5\% ćwiczących (łącznie z piłką siatkową plażową - 7,8\%). Kolejne miejsca pod względem liczby ćwiczących zajmują koszykówka, pływanie i lekkoatletyka (po 4,0\%) oraz karate (łącznie z karate tradycyjnym) - 3,9\%. 

    \subsection{Systemy do zarządzania obiektami sportowymi}
    
        \textbf{Jastrzębik} skupiony jest na oferowaniu użytkownikowi końcowemu obiektów takich jak boiska, korty tenisowe, tory, strzelnice.
        Konkurencyjne systemy to na przykład:
    
        \begin{itemize}
            \item Raz Dwa Projekt -- system rezerwacji kortów tenisowych i squashowych oraz niezbędnego sprzętu z obsługą płatności. Oferuje również pomoc w znalezieniu partnera do gry.
            \\ W porównaniu do naszego przedsięwzięcia jest on mało zróżnicowany pod względem typów obiektów sportowych.
            \item Rezerwuj Sport -- system rezerwacji oraz kupowania karnetów do obiektów sportowych na terenie Wrocławia.
            \\ Głównym minusem tego systemu jest fakt, że dotyczy on obiektów zaledwie w jednym mieście.
            \item Asysto -- oferujący oprogramowanie internetowe pozwalające na uruchomienie rezerwacji online obiektów sportowych na swojej stronie. Jednak nie jest on portalem rezerwacyjnym.
        \end{itemize}
    
\section{Opis firmy i jej działalności}

    Firma \textbf{Jastrzębik} będzie oferować możliwość rezerwacji obiektów sportowo-rekreacyjnych poprzez aplikację internetową. Obiekty te będą one w posiadaniu konkretnych animatorów, którzy mogą wyrazić chęć zarządzania nimi za pomocą niniejszego systemu.
    Użytkownik może zarezerwować taki obiekt w ustalonym czasie za co animator pobiera opłatę.
    Źródłem dochodów firmy będzie prowizja pobierana od animatorów za rezerwacje oraz opłata za wyświetlane za pośrednictwem witryny internetowej reklamy.
    Głównym zadaniem firmy będzie utrzymanie sprawnego działania serwisu internetowego dokonywane poprzez zatrudnionych administratorów. Dodatkowo firma będzie odpowiedzialna za wprowadzanie aktualizacji do wydanego już oprogramowania.

\section{Kontekst funkcjonowania systemu}

    \subsection{Cel} 

        System \textbf{Jastrzębik} to uniwersalna platforma ułatwiająca proces rezerwacji obiektów sportowych na terenie całego kraju. Dzięki dużej bazie obiektów oraz jednym portalu rezerwacyjnym wynajem stanie się dużo prostszy oraz bardziej intuicyjny.
    
        \begin{figure}[H] 
        	\begin{center}
        		\includegraphics[width=0.5\textwidth]{img1/Logo.png}
                \caption{Logo \textbf{Jastrzębika}}
        	\end{center}
        \end{figure}

    \subsection{Zakres}
    
        Zakresem działania systemu jest wynajem obiektów sportowych. Jedna strona, animator, udostępnia obiekt do najmu, druga, użytkownik, dokonuje rezerwacji. Nad poprawnością przeprowadzanych operacji będzie czuwał administrator systemu.

    \subsection{Procesy}
        Wśród najważniejszych procesów znajdą się:
    
        \begin{itemize}
            \item Kreacja oraz udostępnianie swoich obiektów przez animatorów,
            \item Utworzenie rezerwacji oraz administracja nią,
            \item Zgłaszanie nieprawidłowości.
        \end{itemize}

    \subsection{Dokumenty}    
    
        Serwis \textbf{Jastrzębika} zapewni dostęp do następujących dokumentów:
        \begin{itemize}
            \item Regulamin serwisu, zawierający m.in. informację o ochronie danych osobowych,
            \item Regulamin tworzenia i zarządzania obiektami sportowymi,
            \item Regulamin wynajmu obiektów sportowych.
        \end{itemize}
        
        Ponadto, animatorzy mają obowiązek udostępnienia regulaminu każdego swojego obiektu.

\section{Funkcje systemu}

    Wszyscy korzystający z systemu mają możliwość:
    \begin{itemize}
        \item[]
        \begin{itemize}
            \item Logowania i rejestracji,
            \item Podglądu i filtrowania obiektów.
        \end{itemize}
    \end{itemize}
    
    \par Pozostałe funkcje są zależne od roli w systemie według klasyfikacji poniżej:
    \begin{itemize}
        \item Użytkownik
            \begin{itemize}
                \item Przeglądanie i filtrowanie obiektów oraz ich harmonogramu dostępności,
                \item Tworzenie, modyfikacja i anulowanie rezerwacji obiektów,
                \item Kontakt z animatorami oraz administracją.
            \end{itemize}
        \item Animator
            \begin{itemize}
                \item Tworzenie obiektów sportowych w systemie,
                \item Zarządzanie obiektami oraz ich dostępnością,
                \item Usunięcie rezerwacji w przypadku sytuacji kryzysowych.
            \end{itemize}
        \item Administrator
            \begin{itemize}
                \item Zmiana harmonogramu dostępności obiektów,
                \item Blokada konta.
            \end{itemize}
    \end{itemize}


\section{Charakterystyka użytkowników}

    Docelowymi użytkownikami systemu są animator udostępniający obiekty sportowe, użytkownik je rezerwujący oraz administrator. Aby skorzystać z systemu każdy użytkownik musi posiadać aktywne konto.

    Jako \textbf{użytkownika rezerwującego} obiekty sportowe rozumiemy osoby powyżej 16 roku życia zamieszkałe lub przebywające na terenie Polski. Będą oni posiadali ograniczone uprawnienia w systemie - możliwość przeglądania dostępnych kompleksów sportowych, wynajęcie danego obiektu i ewentualne późniejsze edytowanie lub anulowanie rezerwacji. Dodatkowo użytkownik może zgłosić animatora. 

    \textbf{Animatorzy} to użytkownicy udostępniający do wynajmu swoje kompleksy sportowe. Posiadają oni wgląd do własnych oraz możliwość przeglądania wszystkich obiektów. Będą oni posiadali szeroki zakres uprawnień w systemie - przegląd oraz anulowanie rezerwacji powiązanych z ich obiektami, uaktualnianie informacji na temat własnych obiektów oraz ustalanie harmonogramu dostępności. Dodatkowo animator może zgłosić użytkownika rezerwującego. 

    Użytkownikami odpowiedzialnymi za kontrolę nad funkcjonowaniem całego systemu będą \textbf{administratorzy}. Będą oni rozpatrywać wszelkie wykryte nieprawidłowości w produkcie. Otrzymają pełen zakres uprawnień, włączając w to zarządzanie wszystkimi obiektami sportowymi, a także kontami użytkowników. Administrator jako jedyny może zdjąć lub nałożyć blokadę.

\section{Założenia i zależności systemu}

    Aplikacja, aby spełnić wszystkie swoje założenia funkcjonalne i niefunkcjonalne, będzie potrzebowała komunikować się z systemami zewnętrznymi. 

    Jednym z nich będzie serwis internetowy umożliwiający wyszukiwanie obiektów i oglądanie map. Najpopularniejszym tego typu systemem, który prawdopodobnie zostanie wykorzystany w aplikacji, jest serwis Google Maps. Umożliwia on udostępnianie map we własnej aplikacji. Użytkownik będzie mógł sprawdzić położenie danego obiektu - przy jego podglądzie będzie wyświetlana mapa z jego oznaczeniem.

    Drugim z systemów zewnętrznych wykorzystywanych w projektowanej aplikacji będzie serwis umożliwiający dokonanie płatności internetowych. Użytkownik będzie mógł opłacić stworzone przez siebie rezerwacje. Jednym z serwisów dającym możliwość przeprowadzenia płatności, który z dużym prawdopodobieństwem zostanie wykorzystany w aplikacji, jest Paypal - przedsiębiorstwo oferujące usługi płatnicze umożliwiające przedsiębiorcom i konsumentom posiadającym adres e-mail wysyłanie oraz odbieranie płatności przez Internet.

    Oba wymienione systemy zewnętrzne oferują przystępne interfejsy programistyczne.

\section{Wymagania funkcjonalne i niefunkcjonalne}

    Zaimplementowane rozwiązanie powinno spełniać szereg wymagań funkcjonalnych oraz niefunkcjonalnych, aby zapewnić poprawne działanie aplikacji.

    \subsection{Wymagania funkcjonalne} 

        \begin{tabular}{|p{0.45\linewidth}|ccc|}
            \hline
            Wymaganie & Użytkownik & Animator & Administrator \\ \hline \hline
            Rejestracja & \checkmark & \checkmark & \\ \hline
            Logowanie oraz wylogowanie się z systemu & \checkmark & \checkmark & \checkmark \\ \hline
            Utworzenie obiektu sportowego & & \checkmark & \\ \hline
            Podgląd obiektów & \checkmark & \checkmark & \checkmark \\ \hline
            Podgląd własnych obiektów & & \checkmark & \\ \hline
            Zmiana harmonogramu dostępności obiektów & & \checkmark & \checkmark* \\ \hline
            Wgląd do harmonogramu dostępności obiektów & \checkmark & \checkmark & \checkmark \\ \hline
            Rezerwacja wybranego obiektu & \checkmark & & \\ \hline
            Podgląd rezerwacji & \checkmark & \checkmark & \checkmark \\ \hline
            Opłacenie rezerwacji & \checkmark & & \\ \hline
            Modyfikacja rezerwacji & \checkmark & & \checkmark* \\ \hline
            Anulowanie rezerwacji & \checkmark & \checkmark** & \checkmark* \\ \hline
            Dostęp do panelu kontaktu & \checkmark & \checkmark & \checkmark \\ \hline
            Zgłaszanie użytkownika / animatora & \checkmark & \checkmark & \\ \hline
            Blokada konta użytkownika & & & \checkmark \\ \hline
        \end{tabular}
        \\

        \small{* Administrator może zmienić harmonogram oraz modyfikować lub usuwać rezerwacje np. w przypadku wydanego przez Państwo rozporządzenia o zakazie zgromadzeń publicznych lub w innych określonych w regulaminie okolicznościach czy zablokowania użytkowników bądź animatorów z powodu naruszania regulaminu.}
    
        \small{** Animator ma prawo do usunięcia rezerwacji z powodu nieosiągalności obiektu w dowolnym momencie, co zawsze będzie się wiązało z przynajmniej pełnym zwrotem kosztów rezerwacji. }

    \subsection{Wymagania niefunkcjonalne}

        \begin{enumerate}
        	\item Użytkownik nie ma możliwości wglądu w konta innych użytkowników.
        	\item Użytkownik podczas rejestracji musi mieć ukończone co najmniej 16 lat.
        	\item Jedna osoba fizyczna może posiadać tylko jedno konto.
        	\item W przypadku zgłoszeń użytkownika/animatora blokowana jest możliwość rezerwowania/udostępniania obiektów sportowych zgodnie z zasadami określonymi w regulaminie.
        	\item W przypadku dokonania lub odkrycia szkody na terenie obiektu sportowego, użytkownik zobowiązany jest do natychmiastowego zgłoszenia takiego incydentu do animatora.
        	\item W przypadku, gdy podczas użytkowania obiekt sportowy zostanie uszkodzony, rezerwujący zobowiązany jest do pokrycia kosztów naprawy lub wymiany
        \end{enumerate}

\section{Aspekty prawne dotyczące prowadzonej działalności}

    System rezerwacji \textbf{Jastrzębik} udostępni regulamin dla użytkowników oraz animatorów, którego przestrzeganie będzie koniecznie w celu wynajmu i rezerwacji obiektów sportowych. System zostanie stworzony zgodnie z ogólnym rozporządzeniem o ochronie danych, a na stronie zostanie wyświetlony komunikat o zbieranych przez witrynę ciasteczkach. Każda osoba będzie mogła znaleźć dane firmy takie jak nazwa, adres, właściciele, numer NIP czy REGON.

\section{Diagram przypadków użycia}

    \begin{figure}[H] 
    	\begin{center}
    	\includegraphics[width=1\textwidth]{img2/usecase.png}
        \caption{Diagram przypadków użycia}
    	\end{center}
    \end{figure}

\section{ERD}

    \begin{figure}[H] 
    	\begin{center}
    	\includegraphics[width=1\textwidth]{img2/erd.png}
        \caption{Diagram ERD}
    	\end{center}
    \end{figure}

\section{Diagram hierarchii funkcji}

    \begin{figure}[H] 
    	\begin{center}
    		\includegraphics[width=1\textwidth]{img2/HFD.jpg}
            \caption{Diagram hierarchii funkcji}
    	\end{center}
    \end{figure}
    
    \begin{figure}[H] 
    	\begin{center}
    		\includegraphics[width=1\textwidth]{img2/HFD-1.png}
            \caption{Diagram hierarchii funkcji - Zarządzanie kontem}
    	\end{center}
    \end{figure}
    
    \begin{figure}[H] 
    	\begin{center}
    		\includegraphics[width=1\textwidth]{img2/HFD-2.png}
            \caption{Diagram hierarchii funkcji - Zarządzanie obiektami}
    	\end{center}
    \end{figure}
    
    \begin{figure}[H] 
    	\begin{center}
    		\includegraphics[width=1\textwidth]{img2/HFD-3.png}
            \caption{Diagram hierarchii funkcji - zarządzanie rezerwacjami}
    	\end{center}
    \end{figure}
    
    \begin{figure}[H]
    	\begin{center}
    		\includegraphics[width=1\textwidth]{img2/HFD-4.png}
            \caption{Diagram hierarchii funkcji - Zarządzanie użytkownikami}
    	\end{center}
    \end{figure}

\section{Projekt interfejsu}

    Na Rysunku \ref{fig:projekt} zaprezentowany został ogólny projekt interfejsu systemu \textbf{Jastrzębik}:

    \begin{figure}[H]
    	\begin{center}
    		\includegraphics[width=0.8\textwidth]{img2/Projekt.png}
            \caption{Projekt interfejsu serwisu}
            \label{fig:projekt}
    	\end{center}
    \end{figure}

    Będzie się on składać z następujących części:

    \begin{itemize}
        \item Logo -- przycisk nawigacyjny, przekierowujący użytkownika na główną stronę, w postaci obrazka z logo Jastrzębika,
        \item Panel górny - pasek nawigacyjny.
        \item Panel logowania -- w zależności od statusu użytkownika:
            \begin{itemize}
                \item Użytkownik niezalogowany -- moduł wyświetlający panel logowania bądź skrót do strony rejestracyjnej,
                \item Użytkownik zalogowany -- podstawowe informacje o użytkowniku oraz przycisk umożliwiający wylogowanie.
            \end{itemize}
        \item Panel boczny -- filtry ułatwiające wyszukiwane wśród wyświetlonych treści. 
        \item Panel główny -- w zależności od roli użytkownika oraz otwartej aktualnie zakładki:
            \begin{itemize}
                \item Użytkownik niezalogowany -- treściwie przedstawiona oferta i możliwości portalu. 
                \item Klient -- lista obiektów bądź własnych rezerwacji, szczegóły obiektu bądź rezerwacji, panel umożliwiający wyświetlenie bądź zmianę danych konta,
                \item Animator -- lista własnych obiektów, szczegóły danego obiektu i harmonogram, panel umożliwiający wyświetlenie bądź zmianę danych konta.
                \item Administrator -- lista użytkowników systemu, obiektów lub zgłoszeń.
            \end{itemize}
        \item Stopka -- informacje o firmie, dane kontaktowe, skróty.
    \end{itemize}

\section{Technologie}

    Strona serwerowa systemu powstanie w technologii Java (w wersji nie niżej niż Java SE Development Kit 8) z wykorzystaniem szkieletu tworzenia aplikacji Spring, a w szczególności jego rozszerzenia w postaci biblioteki Spring Boot. Zapewni to możliwość uruchomienia aplikacji na różnych systemach operacyjnych, kompatybilność załączanych narzędzi oraz umożliwi sprawne dodawanie nowych funkcjonalności w projektowanej platformie. Zostanie wykorzystany wbudowany w narzędzie Spring kontener aplikacji webowych Apache Tomcat. 
    
    Ze względu na specyfikę danych wykorzystana zostanie relacyjna baza danych, oparta na technologii H2 Database, która zapewnia wysokie wsparcie dla funkcjonalności oferowanych przez Spring Boot.
    
    Zbudowany obraz systemu zostanie umieszczony w kontenerze. Wykorzystana zostanie technologia Docker będąca wiodącym oprogramowaniem służącym do tworzenia, wdrażania i uruchamiania aplikacji rozproszonych.
    
    Strona kliencka oparta będzie na technologiach JavaScript, CSS, HTML - najpopularniejszych rozwiązaniach na rynku, cieszących się dużym wsparciem twórców oraz społeczności wraz z ciągle rosnącym wachlarzem możliwości.

\end{document}
